\documentclass{ximera}
%nieuw 2 september 2026
\title{Vraagstukken elektrisch veld: goniometrische figuren}
\author{Daan Lipkens}

\begin{document}
\begin{abstract}
    Oefeningen op het elektrisch veld met ladingen in rechthoekige driehoeksconfiguraties.
\end{abstract}
\maketitle

%———————————————————————————————————————
% OEF VELD PYTH 1: basis, twee positieve ladingen
\begin{exercise}
    Twee positieve puntladingen $Q_{1} = +3{,}0 \, \mathrm{\mu C}$ en $Q_{2} = +4{,}0 \, \mathrm{\mu C}$ bevinden zich zoals aangegeven op de tekening, met een rechte hoek in punt \textit{P}. De afstand van $Q_1$ tot $P$ is $15 \, \mathrm{cm}$, en de afstand van $Q_2$ tot $P$ is $20 \, \mathrm{cm}$.

    \begin{center}
        \begin{tikzpicture}[scale=1]
            \tkzDefPoint(0,0){P}
            \tkzDefPoint(1.5,0){Q1}
            \tkzDefPoint(0,2.0){Q2}
            \draw[gray,thick] (P) -- (Q1) node[midway,below]{$15\,\mathrm{cm}$};
            \draw[gray,thick] (P) -- (Q2) node[midway,left]{$20\,\mathrm{cm}$};
            \tkzMarkRightAngle[size=0.25](Q1,P,Q2)
            \tkzLabelPoint[below=1.5mm](P){$P$}
            \tkzLabelPoint[below=1.5mm](Q1){$Q_1$}
            \tkzLabelPoint[left=1.5mm](Q2){$Q_2$}
            \draw[charge+] (Q1) circle (\R) node[scale=0.5]{$+$};
            \draw[charge+] (Q2) circle (\R) node[scale=0.5]{$+$};
            \filldraw (P) circle (1.2pt);
        \end{tikzpicture}
    \end{center}

    Bereken de grootte van het totale elektrisch veld in punt \textit{P}.

    \[ E_{\text{res}} = \answer[tolerance=0.05,onlineshowanswerbutton]{1.5} \cdot 10^{6} \, \mathrm{N/C} \]

    \begin{hint}
        Omdat de rechte hoek in $P$ ligt, staan de twee veldvectoren loodrecht op elkaar.
    \end{hint}

    \begin{solution}\nl
        \underline{Gegevens:}
        \begin{align*}
            & Q_{1} = +3{,}0 \, \mathrm{\mu C} = 3{,}0 \cdot 10^{-6} \, \mathrm{C} & & Q_{2} = +4{,}0 \, \mathrm{\mu C} = 4{,}0 \cdot 10^{-6} \, \mathrm{C} \\
            & r_{1} = 15 \, \mathrm{cm} = 0{,}15 \, \mathrm{m} & & r_{2} = 20 \, \mathrm{cm} = 0{,}20 \, \mathrm{m}
        \end{align*}

        \underline{Gevraagd:} $E_{\text{res}}$ in $P$\\

        \underline{Oplossing:}\\
        \begin{align*}
            E_{1} &= k_e \frac{|Q_{1}|}{r_{1}^{2}} = \frac{8{,}99 \cdot 10^{9} \, \mathrm{\frac{N \cdot m^2}{C^2}} \cdot 3{,}0 \cdot 10^{-6} \, \mathrm{C}}{(0{,}15 \, \mathrm{m})^{2}} = 1{,}2 \cdot 10^{6} \, \mathrm{N/C} \\
            E_{2} &= k_e \frac{|Q_{2}|}{r_{2}^{2}} = \frac{8{,}99 \cdot 10^{9} \, \mathrm{\frac{N \cdot m^2}{C^2}} \cdot 4{,}0 \cdot 10^{-6} \, \mathrm{C}}{(0{,}20 \, \mathrm{m})^{2}} = 9{,}0 \cdot 10^{5} \, \mathrm{N/C}
        \end{align*}
        Beide ladingen zijn positief, dus wijzen $\vec{E}_1$ en $\vec{E}_2$ elk weg van hun eigen lading; de rechte hoek in $P$ zorgt ervoor dat ze loodrecht op elkaar staan.
        \begin{align*}
            E_{\text{res}} &= \sqrt{E_1^{2} + E_2^{2}} \simeq 1{,}5 \cdot 10^{6} \, \mathrm{N/C}
        \end{align*}
    \end{solution}
\end{exercise}
%———————————————————————————————————————

%———————————————————————————————————————
% OEF DRIEHOEK 2: velden loodrecht op elkaar (Pythagoras)
\begin{exercise}
    Twee puntladingen $Q_{1} = +3{,}0 \cdot 10^{-6} \, \mathrm{C}$ en $Q_{2} = +6{,}0 \cdot 10^{-6} \, \mathrm{C}$ bevinden zich respectievelijk op de $x$-as en de $y$-as ten opzichte van punt \textit{P}, zoals aangegeven op de tekening. $Q_1$ bevindt zich op $0{,}25 \, \mathrm{m}$ van $P$, en $Q_2$ op $0{,}40 \, \mathrm{m}$ van $P$.

    \begin{center}
        \begin{tikzpicture}[scale=1]
            \tkzDefPoint(0,0){P}
            \tkzDefPoint(2.5,0){Q1}
            \tkzDefPoint(0,4){Q2}
            \draw[gray,thick] (P) -- (Q1);
            \draw[gray,thick] (P) -- (Q2);
            \tkzFillAngle[fill=blue!20,opacity=0.5](Q1,P,Q2)
            \tkzLabelAngle[pos=0.6](Q1,P,Q2){$90^\circ$}
            \tkzLabelPoint[below=1.5mm](P){$P$}
            \tkzLabelPoint[below=1.5mm](Q1){$Q_1$}
            \tkzLabelPoint[left=1.5mm](Q2){$Q_2$}
            \draw[charge+] (Q1) circle (\R) node[scale=0.5]{$+$};
            \draw[charge+] (Q2) circle (\R) node[scale=0.5]{$+$};
            \filldraw (P) circle (1.2pt);
        \end{tikzpicture}
    \end{center}

    Bereken de grootte van het totale elektrisch veld in punt \textit{P}.

    \[ E_{\text{res}} = \answer[tolerance=0.05,onlineshowanswerbutton]{5.5} \cdot 10^{5} \, \mathrm{N/C} \]

    \begin{hint}
        Bereken eerst de grootte van $E_1$ en $E_2$ afzonderlijk met de formule voor het radiaal veld.
    \end{hint}
    \begin{hint}
        Omdat $Q_1$ op de $x$-as en $Q_2$ op de $y$-as ligt (een rechte hoek in $P$), staan $\vec{E}_1$ en $\vec{E}_2$ loodrecht op elkaar. Je kan de resultante dan rechtstreeks berekenen met de stelling van Pythagoras.
    \end{hint}

    \begin{solution}\nl
        \underline{Gegevens:}
        \begin{align*}
            & Q_{1} = +3{,}0 \cdot 10^{-6} \, \mathrm{C} & & Q_{2} = +6{,}0 \cdot 10^{-6} \, \mathrm{C} \\
            & r_{1} = 0{,}25 \, \mathrm{m} & & r_{2} = 0{,}40 \, \mathrm{m}
        \end{align*}

        \underline{Gevraagd:} $E_{\text{res}}$ in $P$\\

        \underline{Oplossing:}\\
        \textbf{Stap 1: Grootte van de afzonderlijke velden}
        \begin{align*}
            E_{1} &= k_e \frac{|Q_{1}|}{r_{1}^{2}} = \frac{8{,}99 \cdot 10^{9} \, \mathrm{\frac{N \cdot m^2}{C^2}} \cdot 3{,}0 \cdot 10^{-6} \, \mathrm{C}}{(0{,}25 \, \mathrm{m})^{2}} = 4{,}3 \cdot 10^{5} \, \mathrm{N/C} \\
            E_{2} &= k_e \frac{|Q_{2}|}{r_{2}^{2}} = \frac{8{,}99 \cdot 10^{9} \, \mathrm{\frac{N \cdot m^2}{C^2}} \cdot 6{,}0 \cdot 10^{-6} \, \mathrm{C}}{(0{,}40 \, \mathrm{m})^{2}} = 3{,}4 \cdot 10^{5} \, \mathrm{N/C}
        \end{align*}
        Beide ladingen zijn positief, dus wijst $\vec{E}_1$ weg van $Q_1$ (langs de $x$-as) en $\vec{E}_2$ weg van $Q_2$ (langs de $y$-as). Deze twee vectoren staan dus loodrecht op elkaar.

        \textbf{Stap 2: Stelling van Pythagoras toepassen}
        \begin{align*}
            E_{\text{res}} &= \sqrt{E_1^{2} + E_2^{2}} \\
                            &= \sqrt{(4{,}3 \cdot 10^{5} \, \mathrm{N/C})^{2} + (3{,}4 \cdot 10^{5} \, \mathrm{N/C})^{2}} \\
                            &\simeq 5{,}5 \cdot 10^{5} \, \mathrm{N/C}
        \end{align*}
    \end{solution}
\end{exercise}
%———————————————————————————————————————

%———————————————————————————————————————
%  OEF DRIEHOEK PYTH 3: eenvoudigste geval, 2 positieve ladingen
\begin{exercise}
    Twee positieve puntladingen $Q_{1} = +2{,}0 \cdot 10^{-6} \, \mathrm{C}$ en $Q_{2} = +5{,}0 \cdot 10^{-6} \, \mathrm{C}$ bevinden zich zoals aangegeven op de tekening, met een rechte hoek in punt \textit{P}. De afstand van $Q_1$ tot $P$ is $0{,}20 \, \mathrm{m}$, en de afstand van $Q_2$ tot $P$ is $0{,}30 \, \mathrm{m}$.
 
    \begin{center}
        \begin{tikzpicture}[scale=1]
            \tkzDefPoint(0,0){P}
            \tkzDefPoint(2.0,0){Q1}
            \tkzDefPoint(0,3.0){Q2}
            \draw[gray,thick] (P) -- (Q1) node[midway,below]{$0{,}20\,\mathrm{m}$};
            \draw[gray,thick] (P) -- (Q2) node[midway,left]{$0{,}30\,\mathrm{m}$};
            \tkzMarkRightAngle[size=0.3](Q1,P,Q2)
            \tkzLabelPoint[below=1.5mm](P){$P$}
            \tkzLabelPoint[below=1.5mm](Q1){$Q_1$}
            \tkzLabelPoint[left=1.5mm](Q2){$Q_2$}
            \draw[charge+] (Q1) circle (\R) node[scale=0.5]{$+$};
            \draw[charge+] (Q2) circle (\R) node[scale=0.5]{$+$};
            \filldraw (P) circle (1.2pt);
        \end{tikzpicture}
    \end{center}
 
    Bereken de grootte van het totale elektrisch veld in punt \textit{P}.
 
    \[ E_{\text{res}} = \answer[tolerance=0.1,onlineshowanswerbutton]{6.7} \cdot 10^{5} \, \mathrm{N/C} \]

 
    \begin{solution}\nl
        \underline{Gegevens:}
        \begin{align*}
            & Q_{1} = +2{,}0 \cdot 10^{-6} \, \mathrm{C} & & Q_{2} = +5{,}0 \cdot 10^{-6} \, \mathrm{C} \\
            & r_{1} = 0{,}20 \, \mathrm{m} & & r_{2} = 0{,}30 \, \mathrm{m}
        \end{align*}
 
        \underline{Gevraagd:} $E_{\text{res}}$ in $P$\\
 
        \underline{Oplossing:}\\
        \textbf{Stap 1: Grootte van de afzonderlijke velden}
        \begin{align*}
            E_{1} &= k_e \frac{|Q_{1}|}{r_{1}^{2}} = \frac{8{,}99 \cdot 10^{9} \, \mathrm{\frac{N \cdot m^2}{C^2}} \cdot 2{,}0 \cdot 10^{-6} \, \mathrm{C}}{(0{,}20 \, \mathrm{m})^{2}} = 4{,}5 \cdot 10^{5} \, \mathrm{N/C} \\
            E_{2} &= k_e \frac{|Q_{2}|}{r_{2}^{2}} = \frac{8{,}99 \cdot 10^{9} \, \mathrm{\frac{N \cdot m^2}{C^2}} \cdot 5{,}0 \cdot 10^{-6} \, \mathrm{C}}{(0{,}30 \, \mathrm{m})^{2}} = 5{,}0 \cdot 10^{5} \, \mathrm{N/C}
        \end{align*}
        Beide ladingen zijn positief, dus wijst $\vec{E}_1$ weg van $Q_1$ en $\vec{E}_2$ weg van $Q_2$. Omdat de rechte hoek in $P$ ligt, staan deze twee vectoren loodrecht op elkaar.
 
        \textbf{Stap 2: Stelling van Pythagoras toepassen}
        \begin{align*}
            E_{\text{res}} &= \sqrt{E_1^{2} + E_2^{2}} \\
                            &= \sqrt{(4{,}5 \cdot 10^{5} \, \mathrm{N/C})^{2} + (5{,}0 \cdot 10^{5} \, \mathrm{N/C})^{2}} \\
                            &\simeq 6{,}7 \cdot 10^{5} \, \mathrm{N/C}
        \end{align*}
    \end{solution}
\end{exercise}
%———————————————————————————————————————
 
%———————————————————————————————————————
% OEF DRIEHOEK PYTH 4: rechte hoek in P, gemengde tekens
\begin{exercise}
    Een positieve puntlading $Q_{1} = +3{,}0 \cdot 10^{-6} \, \mathrm{C}$ en een negatieve puntlading $Q_{3} = -4{,}0 \cdot 10^{-6} \, \mathrm{C}$ bevinden zich zoals aangegeven op de tekening, met een rechte hoek in punt \textit{P}. De afstand van $Q_1$ tot $P$ is $0{,}15 \, \mathrm{m}$, en de afstand van $Q_3$ tot $P$ is $0{,}25 \, \mathrm{m}$.
 
    \begin{center}
        \begin{tikzpicture}[scale=1.4]
            \tkzDefPoint(0,0){P}
            \tkzDefPoint(1.5,0){Q1}
            \tkzDefPoint(0,2.5){Q3}
            \draw[gray,thick] (P) -- (Q1) node[midway,below]{$0{,}15\,\mathrm{m}$};
            \draw[gray,thick] (P) -- (Q3) node[midway,left]{$0{,}25\,\mathrm{m}$};
            \tkzMarkRightAngle[size=0.25](Q1,P,Q3)
            \tkzLabelPoint[below=1.5mm](P){$P$}
            \tkzLabelPoint[below=1.5mm](Q1){$Q_1$}
            \tkzLabelPoint[left=1.5mm](Q3){$Q_3$}
            \draw[charge+] (Q1) circle (\R) node[scale=0.5]{$+$};
            \draw[charge-] (Q3) circle (\R) node[scale=0.5]{$-$};
            \filldraw (P) circle (1.2pt);
        \end{tikzpicture}
    \end{center}
 
    \begin{question}
        Maak een schets van $\vec{E}_1$ en $\vec{E}_3$ in punt $P$. Welke bewering klopt?
 
        \begin{multipleChoice}
            \choice{$\vec{E}_1$ en $\vec{E}_3$ wijzen beide van $P$ weg, in dezelfde zin.}
            \choice[correct]{$\vec{E}_1$ wijst van $P$ weg (langs de lijn naar $Q_1$), $\vec{E}_3$ wijst van $P$ naar $Q_3$ toe. Ze staan nog steeds loodrecht op elkaar.}
            \choice{Omdat de ladingen een verschillend teken hebben, kan je de stelling van Pythagoras hier niet gebruiken.}
        \end{multipleChoice}
 
        \begin{hint}
            Kijk naar het teken van elke lading afzonderlijk om de zin van elk veld te bepalen. De hoek tussen de twee \textit{lijnstukken} $PQ_1$ en $PQ_3$ verandert niet door het teken van de ladingen.
        \end{hint}
    \end{question}
 
    \begin{question}
        Bereken de grootte van het totale elektrisch veld in punt \textit{P} als daar een lading $Q_2 = +2{,}0 \cdot 10^{-6} \, \mathrm{C}$ is geplaatst?
 
        \[ E_{\text{res}} = \answer[tolerance=0.05,onlineshowanswerbutton]{1.3} \cdot 10^{6} \, \mathrm{N/C} \]
 
        \begin{hint}
            $Q_2$ staat in punt \textit{P} en heeft dus geen invloed op het elektrisch veld in punt \textit{P} zelf. Deze lading heb je dus niet nodig.
        \end{hint}
 
        \begin{solution}\nl
            \underline{Gegevens:}
            \begin{align*}
                & Q_{1} = +3{,}0 \cdot 10^{-6} \, \mathrm{C} & & Q_{3} = -4{,}0 \cdot 10^{-6} \, \mathrm{C} \\
                & r_{1} = 0{,}15 \, \mathrm{m} & & r_{3} = 0{,}25 \, \mathrm{m}
            \end{align*}
 
            \underline{Gevraagd:} $E_{\text{res}}$ in $P$\\
 
            \underline{Oplossing:}\\
            \textbf{Stap 1: Grootte van de afzonderlijke velden}
            \begin{align*}
                E_{1} &= k_e \frac{|Q_{1}|}{r_{1}^{2}} = \frac{8{,}99 \cdot 10^{9} \, \mathrm{\frac{N \cdot m^2}{C^2}} \cdot 3{,}0 \cdot 10^{-6} \, \mathrm{C}}{(0{,}15 \, \mathrm{m})^{2}} = 1{,}2 \cdot 10^{6} \, \mathrm{N/C} \\
                E_{3} &= k_e \frac{|Q_{3}|}{r_{3}^{2}} = \frac{8{,}99 \cdot 10^{9} \, \mathrm{\frac{N \cdot m^2}{C^2}} \cdot 4{,}0 \cdot 10^{-6} \, \mathrm{C}}{(0{,}25 \, \mathrm{m})^{2}} = 5{,}8 \cdot 10^{5} \, \mathrm{N/C}
            \end{align*}
            $Q_1$ is positief: $\vec{E}_1$ wijst van $P$ weg (langs $PQ_1$). $Q_3$ is negatief: $\vec{E}_3$ wijst van $P$ naar $Q_3$ toe (langs $PQ_3$). Beide lijnstukken $PQ_1$ en $PQ_3$ staan loodrecht op elkaar, dus staan ook $\vec{E}_1$ en $\vec{E}_3$ loodrecht op elkaar.
 
            \textbf{Stap 2: Stelling van Pythagoras toepassen}
            \begin{align*}
                E_{\text{res}} &= \sqrt{E_1^{2} + E_3^{2}} \\
                                &= \sqrt{(1{,}2 \cdot 10^{6} \, \mathrm{N/C})^{2} + (5{,}8 \cdot 10^{5} \, \mathrm{N/C})^{2}} \\
                                &\simeq 1{,}3 \cdot 10^{6} \, \mathrm{N/C}
            \end{align*}
        \end{solution}
    \end{question}
\end{exercise}
%———————————————————————————————————————
 
%———————————————————————————————————————
% OEF DRIEHOEK PYTH 5: rechte hoek in P, ontbrekende afstand via Pythagoras
\begin{exercise}
    Twee puntladingen $Q_{1} = +6{,}0 \cdot 10^{-6} \, \mathrm{C}$ en $Q_{2} = -8{,}0 \cdot 10^{-6} \, \mathrm{C}$ bevinden zich zoals aangegeven op de tekening, met een rechte hoek in punt \textit{P}. De afstand van $Q_1$ tot $P$ is $0{,}30 \, \mathrm{m}$. De afstand tussen $Q_1$ en $Q_2$ zelf (dus niet tot $P$) bedraagt $0{,}50 \, \mathrm{m}$.
 
    \begin{center}
        \begin{tikzpicture}[scale=1.2]
            \tkzDefPoint(0,0){P}
            \tkzDefPoint(3.0,0){Q1}
            \tkzDefPoint(0,4.0){Q2}
            \draw[gray,thick] (P) -- (Q1) node[midway,below]{$0{,}30\,\mathrm{m}$};
            \draw[gray,dashed] (P) -- (Q2) node[midway,left]{$?$};
            \draw[gray,thick] (Q1) -- (Q2) node[midway,above right]{$0{,}50\,\mathrm{m}$};
            \tkzMarkRightAngle[size=0.3](Q1,P,Q2)
            \tkzLabelPoint[below=1.5mm](P){$P$}
            \tkzLabelPoint[below=1.5mm](Q1){$Q_1$}
            \tkzLabelPoint[left=1.5mm](Q2){$Q_2$}
            \draw[charge+] (Q1) circle (\R) node[scale=0.5]{$+$};
            \draw[charge-] (Q2) circle (\R) node[scale=0.5]{$-$};
            \filldraw (P) circle (1.2pt);
        \end{tikzpicture}
    \end{center}
 
    \begin{question}
        Bereken eerst de afstand van $Q_2$ tot $P$.
 
        \[ r_{2} = \answer[tolerance=0.01,onlineshowanswerbutton]{0.40} \, \mathrm{m} \]
 
        \begin{hint}
            De driehoek gevormd door $P$, $Q_1$ en $Q_2$ is rechthoekig, met de rechte hoek in $P$. De gegeven afstand van $0{,}50 \, \mathrm{m}$ tussen $Q_1$ en $Q_2$ is dus de schuine zijde van die driehoek.
        \end{hint}
 
        \begin{solution}\nl
            De driehoek $PQ_1Q_2$ is rechthoekig met de rechte hoek in $P$. De zijde $Q_1Q_2 = 0{,}50 \, \mathrm{m}$ is dan de schuine zijde. Met de stelling van Pythagoras:
            \begin{align*}
                r_{2} &= \sqrt{(Q_1Q_2)^{2} - r_{1}^{2}} \\
                      &= \sqrt{(0{,}50 \, \mathrm{m})^{2} - (0{,}30 \, \mathrm{m})^{2}} \\
                      &= 0{,}40 \, \mathrm{m}
            \end{align*}
        \end{solution}
    \end{question}
 
    \begin{question}
        Bereken nu de grootte van het totale elektrisch veld in punt \textit{P}.
 
        \[ E_{\text{res}} = \answer[tolerance=0.05,onlineshowanswerbutton]{7.5} \cdot 10^{5} \, \mathrm{N/C} \]
 
        \begin{hint}
            Bereken $E_1$ en $E_2$ met de afstanden $r_1 = 0{,}30 \, \mathrm{m}$ en $r_2 = 0{,}40 \, \mathrm{m}$. Omdat de rechte hoek zich in $P$ bevindt, staan de twee veldvectoren loodrecht op elkaar.
        \end{hint}
 
        \begin{solution}\nl
            \underline{Gegevens:}
            \begin{align*}
                & Q_{1} = +6{,}0 \cdot 10^{-6} \, \mathrm{C} & & Q_{2} = -8{,}0 \cdot 10^{-6} \, \mathrm{C} \\
                & r_{1} = 0{,}30 \, \mathrm{m} & & r_{2} = 0{,}40 \, \mathrm{m}
            \end{align*}
 
            \underline{Gevraagd:} $E_{\text{res}}$ in $P$\\
 
            \underline{Oplossing:}\\
            \textbf{Stap 1: Grootte van de afzonderlijke velden}
            \begin{align*}
                E_{1} &= k_e \frac{|Q_{1}|}{r_{1}^{2}} = \frac{8{,}99 \cdot 10^{9} \, \mathrm{\frac{N \cdot m^2}{C^2}} \cdot 6{,}0 \cdot 10^{-6} \, \mathrm{C}}{(0{,}30 \, \mathrm{m})^{2}} = 6{,}0 \cdot 10^{5} \, \mathrm{N/C} \\
                E_{2} &= k_e \frac{|Q_{2}|}{r_{2}^{2}} = \frac{8{,}99 \cdot 10^{9} \, \mathrm{\frac{N \cdot m^2}{C^2}} \cdot 8{,}0 \cdot 10^{-6} \, \mathrm{C}}{(0{,}40 \, \mathrm{m})^{2}} = 4{,}5 \cdot 10^{5} \, \mathrm{N/C}
            \end{align*}
            $Q_1$ is positief: $\vec{E}_1$ wijst van $P$ weg. $Q_2$ is negatief: $\vec{E}_2$ wijst van $P$ naar $Q_2$ toe. Omdat de rechte hoek in $P$ ligt, staan beide vectoren loodrecht op elkaar.
 
            \textbf{Stap 2: Stelling van Pythagoras toepassen}
            \begin{align*}
                E_{\text{res}} &= \sqrt{E_1^{2} + E_2^{2}} \\
                                &= \sqrt{(6{,}0 \cdot 10^{5} \, \mathrm{N/C})^{2} + (4{,}5 \cdot 10^{5} \, \mathrm{N/C})^{2}} \\
                                &\simeq 7{,}5 \cdot 10^{5} \, \mathrm{N/C}
            \end{align*}
        \end{solution}
    \end{question}
\end{exercise}
%———————————————————————————————————————

%———————————————————————————————————————
% OEF VELD PYTH 6: basis, twee positieve ladingen
\begin{exercise}
    Twee positieve puntladingen $Q_{1} = +3{,}0 \, \mathrm{\mu C}$ en $Q_{2} = +4{,}0 \, \mathrm{\mu C}$ bevinden zich zoals aangegeven op de tekening, met een rechte hoek in punt \textit{P}. De afstand van $Q_1$ tot $P$ is $15 \, \mathrm{cm}$, en de afstand van $Q_2$ tot $P$ is $20 \, \mathrm{cm}$.

    \begin{center}
        \begin{tikzpicture}[scale=1]
            \tkzDefPoint(0,0){P}
            \tkzDefPoint(1.5,0){Q1}
            \tkzDefPoint(0,2.0){Q2}
            \draw[gray,thick] (P) -- (Q1) node[midway,below]{$15\,\mathrm{cm}$};
            \draw[gray,thick] (P) -- (Q2) node[midway,left]{$20\,\mathrm{cm}$};
            \tkzMarkRightAngle[size=0.25](Q1,P,Q2)
            \tkzLabelPoint[below=1.5mm](P){$P$}
            \tkzLabelPoint[below=1.5mm](Q1){$Q_1$}
            \tkzLabelPoint[left=1.5mm](Q2){$Q_2$}
            \draw[charge+] (Q1) circle (\R) node[scale=0.5]{$+$};
            \draw[charge+] (Q2) circle (\R) node[scale=0.5]{$+$};
            \filldraw (P) circle (1.2pt);
        \end{tikzpicture}
    \end{center}

    Bereken de grootte van het totale elektrisch veld in punt \textit{P}.

    \[ E_{\text{res}} = \answer[tolerance=0.05,onlineshowanswerbutton]{1.5} \cdot 10^{6} \, \mathrm{N/C} \]

    \begin{hint}
        Omdat de rechte hoek in $P$ ligt, staan de twee veldvectoren loodrecht op elkaar.
    \end{hint}

    \begin{solution}\nl
        \underline{Gegevens:}
        \begin{align*}
            & Q_{1} = +3{,}0 \, \mathrm{\mu C} = 3{,}0 \cdot 10^{-6} \, \mathrm{C} & & Q_{2} = +4{,}0 \, \mathrm{\mu C} = 4{,}0 \cdot 10^{-6} \, \mathrm{C} \\
            & r_{1} = 15 \, \mathrm{cm} = 0{,}15 \, \mathrm{m} & & r_{2} = 20 \, \mathrm{cm} = 0{,}20 \, \mathrm{m}
        \end{align*}

        \underline{Gevraagd:} $E_{\text{res}}$ in $P$\\

        \underline{Oplossing:}\\
        \begin{align*}
            E_{1} &= k_e \frac{|Q_{1}|}{r_{1}^{2}} = \frac{8{,}99 \cdot 10^{9} \, \mathrm{\frac{N \cdot m^2}{C^2}} \cdot 3{,}0 \cdot 10^{-6} \, \mathrm{C}}{(0{,}15 \, \mathrm{m})^{2}} = 1{,}2 \cdot 10^{6} \, \mathrm{N/C} \\
            E_{2} &= k_e \frac{|Q_{2}|}{r_{2}^{2}} = \frac{8{,}99 \cdot 10^{9} \, \mathrm{\frac{N \cdot m^2}{C^2}} \cdot 4{,}0 \cdot 10^{-6} \, \mathrm{C}}{(0{,}20 \, \mathrm{m})^{2}} = 9{,}0 \cdot 10^{5} \, \mathrm{N/C}
        \end{align*}
        Beide ladingen zijn positief, dus wijzen $\vec{E}_1$ en $\vec{E}_2$ elk weg van hun eigen lading; de rechte hoek in $P$ zorgt ervoor dat ze loodrecht op elkaar staan.
        \begin{align*}
            E_{\text{res}} &= \sqrt{E_1^{2} + E_2^{2}} \simeq 1{,}5 \cdot 10^{6} \, \mathrm{N/C}
        \end{align*}
    \end{solution}
\end{exercise}
%———————————————————————————————————————

%———————————————————————————————————————
% OEF VELD PYTH 2: tegengesteld teken, nC en mm-schaal
\begin{exercise}
    Een positieve puntlading $Q_{1} = +800 \, \mathrm{nC}$ en een negatieve puntlading $Q_{3} = -500 \, \mathrm{nC}$ bevinden zich zoals aangegeven op de tekening, met een rechte hoek in punt \textit{P}. De afstand van $Q_1$ tot $P$ is $80 \, \mathrm{mm}$, en de afstand van $Q_3$ tot $P$ is $12 \, \mathrm{cm}$.

    \begin{center}
        \begin{tikzpicture}[scale=1.3]
            \tkzDefPoint(0,0){P}
            \tkzDefPoint(1.6,0){Q1}
            \tkzDefPoint(0,2.4){Q3}
            \draw[gray,thick] (P) -- (Q1) node[midway,below]{$80\,\mathrm{mm}$};
            \draw[gray,thick] (P) -- (Q3) node[midway,left]{$12\,\mathrm{cm}$};
            \tkzMarkRightAngle[size=0.2](Q1,P,Q3)
            \tkzLabelPoint[below=1.5mm](P){$P$}
            \tkzLabelPoint[below=1.5mm](Q1){$Q_1$}
            \tkzLabelPoint[left=1.5mm](Q3){$Q_3$}
            \draw[charge+] (Q1) circle (\R) node[scale=0.5]{$+$};
            \draw[charge-] (Q3) circle (\R) node[scale=0.5]{$-$};
            \filldraw (P) circle (1.2pt);
        \end{tikzpicture}
    \end{center}

    Bereken de grootte van het totale elektrisch veld in punt \textit{P}.

    \[ E_{\text{res}} = \answer[tolerance=0.05,onlineshowanswerbutton]{1.2} \cdot 10^{6} \, \mathrm{N/C} \]

    \begin{hint}
        Ondanks het verschil in teken staan $\vec{E}_1$ en $\vec{E}_3$ nog steeds loodrecht op elkaar, omdat de rechte hoek in $P$ ligt.
    \end{hint}

    \begin{solution}\nl
        \underline{Gegevens:}
        \begin{align*}
            & Q_{1} = +800 \, \mathrm{nC} = 8{,}0 \cdot 10^{-7} \, \mathrm{C} & & Q_{3} = -500 \, \mathrm{nC} = 5{,}0 \cdot 10^{-7} \, \mathrm{C} \\
            & r_{1} = 80 \, \mathrm{mm} = 0{,}080 \, \mathrm{m} & & r_{3} = 12 \, \mathrm{cm} = 0{,}12 \, \mathrm{m}
        \end{align*}

        \underline{Gevraagd:} $E_{\text{res}}$ in $P$\\

        \underline{Oplossing:}\\
        \begin{align*}
            E_{1} &= k_e \frac{|Q_{1}|}{r_{1}^{2}} = \frac{8{,}99 \cdot 10^{9} \, \mathrm{\frac{N \cdot m^2}{C^2}} \cdot 8{,}0 \cdot 10^{-7} \, \mathrm{C}}{(0{,}080 \, \mathrm{m})^{2}} = 1{,}1 \cdot 10^{6} \, \mathrm{N/C} \\
            E_{3} &= k_e \frac{|Q_{3}|}{r_{3}^{2}} = \frac{8{,}99 \cdot 10^{9} \, \mathrm{\frac{N \cdot m^2}{C^2}} \cdot 5{,}0 \cdot 10^{-7} \, \mathrm{C}}{(0{,}12 \, \mathrm{m})^{2}} = 3{,}1 \cdot 10^{5} \, \mathrm{N/C}
        \end{align*}
        $\vec{E}_1$ wijst weg van $Q_1$, $\vec{E}_3$ wijst naar $Q_3$ toe — de rechte hoek in $P$ zorgt er wel voor dat ze loodrecht op elkaar blijven staan.
        \begin{align*}
            E_{\text{res}} &= \sqrt{E_1^{2} + E_3^{2}} \simeq 1{,}2 \cdot 10^{6} \, \mathrm{N/C}
        \end{align*}
    \end{solution}
\end{exercise}
%———————————————————————————————————————

%———————————————————————————————————————
% OEF VELD PYTH 3: rechthoek, P in een hoekpunt
\begin{exercise}
    Twee puntladingen $Q_{A} = +5{,}0 \, \mathrm{\mu C}$ en $Q_{B} = -6{,}0 \, \mathrm{\mu C}$ bevinden zich in twee hoekpunten van een rechthoek, aangrenzend aan hoekpunt \textit{P}. De zijde $PQ_A$ is $18 \, \mathrm{cm}$ lang, de zijde $PQ_B$ is $24 \, \mathrm{cm}$ lang.

    \begin{center}
        \begin{tikzpicture}[scale=0.7]
            \tkzDefPoint(0,0){P}
            \tkzDefPoint(3.6,0){A}
            \tkzDefPoint(0,4.8){B}
            \tkzDefPoint(3.6,4.8){C}
            \tkzDrawPolygon(P,A,C,B)
            \tkzMarkRightAngle[size=0.3](A,P,B)
            \tkzLabelPoint[below=1.5mm](P){$P$}
            \tkzLabelPoint[below=1.5mm](A){$Q_A$}
            \tkzLabelPoint[left=1.5mm](B){$Q_B$}
            \draw[charge+] (A) circle (\R) node[scale=0.5]{$+$};
            \draw[charge-] (B) circle (\R) node[scale=0.5]{$-$};
            \filldraw (P) circle (1.5pt);
        \end{tikzpicture}
    \end{center}

    Bereken de grootte van het totale elektrisch veld in punt \textit{P}, veroorzaakt door $Q_A$ en $Q_B$.

    \[ E_{\text{res}} = \answer[tolerance=0.05,onlineshowanswerbutton]{1.7} \cdot 10^{6} \, \mathrm{N/C} \]

    \begin{hint}
        De zijden van een rechthoek staan loodrecht op elkaar, dus doet zich hier hetzelfde voor als in de vorige oefeningen.
    \end{hint}

    \begin{solution}\nl
        \underline{Gegevens:}
        \begin{align*}
            & Q_{A} = +5{,}0 \, \mathrm{\mu C} = 5{,}0 \cdot 10^{-6} \, \mathrm{C} & & Q_{B} = -6{,}0 \, \mathrm{\mu C} = 6{,}0 \cdot 10^{-6} \, \mathrm{C} \\
            & r_{A} = 18 \, \mathrm{cm} = 0{,}18 \, \mathrm{m} & & r_{B} = 24 \, \mathrm{cm} = 0{,}24 \, \mathrm{m}
        \end{align*}

        \underline{Gevraagd:} $E_{\text{res}}$ in $P$\\

        \underline{Oplossing:}\\
        \begin{align*}
            E_{A} &= k_e \frac{|Q_{A}|}{r_{A}^{2}} = \frac{8{,}99 \cdot 10^{9} \, \mathrm{\frac{N \cdot m^2}{C^2}} \cdot 5{,}0 \cdot 10^{-6} \, \mathrm{C}}{(0{,}18 \, \mathrm{m})^{2}} = 1{,}4 \cdot 10^{6} \, \mathrm{N/C} \\
            E_{B} &= k_e \frac{|Q_{B}|}{r_{B}^{2}} = \frac{8{,}99 \cdot 10^{9} \, \mathrm{\frac{N \cdot m^2}{C^2}} \cdot 6{,}0 \cdot 10^{-6} \, \mathrm{C}}{(0{,}24 \, \mathrm{m})^{2}} = 9{,}4 \cdot 10^{5} \, \mathrm{N/C}
        \end{align*}
        De zijden $PQ_A$ en $PQ_B$ van de rechthoek staan loodrecht op elkaar, dus ook $\vec{E}_A$ en $\vec{E}_B$.
        \begin{align*}
            E_{\text{res}} &= \sqrt{E_A^{2} + E_B^{2}} \simeq 1{,}7 \cdot 10^{6} \, \mathrm{N/C}
        \end{align*}
    \end{solution}
\end{exercise}
%———————————————————————————————————————

%———————————————————————————————————————
% OEF VELD PYTH 4: vierkant, gelijke afstanden
\begin{exercise}
    Twee puntladingen $Q_{1} = +2{,}0 \, \mathrm{\mu C}$ en $Q_{2} = -9{,}0 \, \mathrm{\mu C}$ bevinden zich in twee hoekpunten van een vierkant met zijde $25 \, \mathrm{cm}$, aangrenzend aan hoekpunt \textit{P}.

    \begin{center}
        \begin{tikzpicture}[scale=0.5]
            \tkzDefPoint(0,0){P}
            \tkzDefPoint(5,0){Q1}
            \tkzDefPoint(0,5){Q2}
            \tkzDefPoint(5,5){C}
            \tkzDrawPolygon(P,Q1,C,Q2)
            \tkzMarkRightAngle[size=0.4](Q1,P,Q2)
            \tkzLabelPoint[below=1.5mm](P){$P$}
            \tkzLabelPoint[below=1.5mm](Q1){$Q_1$}
            \tkzLabelPoint[left=1.5mm](Q2){$Q_2$}
            \draw[charge+] (Q1) circle (\R) node[scale=0.5]{$+$};
            \draw[charge-] (Q2) circle (\R) node[scale=0.5]{$-$};
            \filldraw (P) circle (1.5pt);
        \end{tikzpicture}
    \end{center}

    Bereken de grootte van het totale elektrisch veld in punt \textit{P}.

    \[ E_{\text{res}} = \answer[tolerance=0.05,onlineshowanswerbutton]{1.3} \cdot 10^{6} \, \mathrm{N/C} \]

    \begin{hint}
        Bij een vierkant zijn beide afstanden tot $P$ gelijk, maar de ladingen zelf zijn dat niet: reken beide velden apart uit.
    \end{hint}

    \begin{solution}\nl
        \underline{Gegevens:}
        \begin{align*}
            & Q_{1} = +2{,}0 \, \mathrm{\mu C} = 2{,}0 \cdot 10^{-6} \, \mathrm{C} & & Q_{2} = -9{,}0 \, \mathrm{\mu C} = 9{,}0 \cdot 10^{-6} \, \mathrm{C} \\
            & r_{1} = r_{2} = 25 \, \mathrm{cm} = 0{,}25 \, \mathrm{m} & &
        \end{align*}

        \underline{Gevraagd:} $E_{\text{res}}$ in $P$\\

        \underline{Oplossing:}\\
        \begin{align*}
            E_{1} &= k_e \frac{|Q_{1}|}{r_{1}^{2}} = \frac{8{,}99 \cdot 10^{9} \, \mathrm{\frac{N \cdot m^2}{C^2}} \cdot 2{,}0 \cdot 10^{-6} \, \mathrm{C}}{(0{,}25 \, \mathrm{m})^{2}} = 2{,}9 \cdot 10^{5} \, \mathrm{N/C} \\
            E_{2} &= k_e \frac{|Q_{2}|}{r_{2}^{2}} = \frac{8{,}99 \cdot 10^{9} \, \mathrm{\frac{N \cdot m^2}{C^2}} \cdot 9{,}0 \cdot 10^{-6} \, \mathrm{C}}{(0{,}25 \, \mathrm{m})^{2}} = 1{,}3 \cdot 10^{6} \, \mathrm{N/C}
        \end{align*}
        De zijden van het vierkant die in $P$ samenkomen, staan loodrecht op elkaar, dus ook $\vec{E}_1$ en $\vec{E}_2$.
        \begin{align*}
            E_{\text{res}} &= \sqrt{E_1^{2} + E_2^{2}} \simeq 1{,}3 \cdot 10^{6} \, \mathrm{N/C}
        \end{align*}
    \end{solution}
\end{exercise}
%———————————————————————————————————————

%———————————————————————————————————————
% OEF VELD PYTH 5: ontbrekende afstand via Pythagoras, meerdelig
\begin{exercise}
    Twee puntladingen $Q_{1} = +7{,}0 \, \mathrm{\mu C}$ en $Q_{2} = -25 \, \mathrm{\mu C}$ bevinden zich zoals aangegeven op de tekening, met een rechte hoek in punt \textit{P}. De afstand van $Q_1$ tot $P$ is $18 \, \mathrm{cm}$. De afstand tussen $Q_1$ en $Q_2$ zelf (dus niet tot $P$) bedraagt $30 \, \mathrm{cm}$.

    \begin{center}
        \begin{tikzpicture}[scale=1]
            \tkzDefPoint(0,0){P}
            \tkzDefPoint(1.8,0){Q1}
            \tkzDefPoint(0,2.4){Q2}
            \draw[gray,thick] (P) -- (Q1) node[midway,below]{$18\,\mathrm{cm}$};
            \draw[gray,dashed] (P) -- (Q2) node[midway,left]{$?$};
            \draw[gray,thick] (Q1) -- (Q2) node[midway,above right]{$30\,\mathrm{cm}$};
            \tkzMarkRightAngle[size=0.25](Q1,P,Q2)
            \tkzLabelPoint[below=1.5mm](P){$P$}
            \tkzLabelPoint[below=1.5mm](Q1){$Q_1$}
            \tkzLabelPoint[left=1.5mm](Q2){$Q_2$}
            \draw[charge+] (Q1) circle (\R) node[scale=0.5]{$+$};
            \draw[charge-] (Q2) circle (\R) node[scale=0.5]{$-$};
            \filldraw (P) circle (1.2pt);
        \end{tikzpicture}
    \end{center}

    \begin{question}
        Bereken eerst de afstand van $Q_2$ tot $P$.

        \[ r_{2} = \answer[tolerance=0.01,onlineshowanswerbutton]{0.24} \, \mathrm{m} \]

        \begin{hint}
            De gegeven afstand van $30 \, \mathrm{cm}$ tussen $Q_1$ en $Q_2$ is de schuine zijde van de rechthoekige driehoek $PQ_1Q_2$.
        \end{hint}

        \begin{solution}\nl
            \begin{align*}
                r_{2} &= \sqrt{(0{,}30 \, \mathrm{m})^{2} - (0{,}18 \, \mathrm{m})^{2}} = 0{,}24 \, \mathrm{m}
            \end{align*}
        \end{solution}
    \end{question}

    \begin{question}
        Bereken nu de grootte van het totale elektrisch veld in punt \textit{P}.

        \[ E_{\text{res}} = \answer[tolerance=0.1,onlineshowanswerbutton]{4.4} \cdot 10^{6} \, \mathrm{N/C} \]

        \begin{solution}\nl
            \underline{Gegevens:}
            \begin{align*}
                & Q_{1} = +7{,}0 \, \mathrm{\mu C} = 7{,}0 \cdot 10^{-6} \, \mathrm{C} & & Q_{2} = -25 \, \mathrm{\mu C} = 2{,}5 \cdot 10^{-5} \, \mathrm{C} \\
                & r_{1} = 0{,}18 \, \mathrm{m} & & r_{2} = 0{,}24 \, \mathrm{m}
            \end{align*}

            \underline{Gevraagd:} $E_{\text{res}}$ in $P$\\

            \underline{Oplossing:}\\
            \begin{align*}
                E_{1} &= k_e \frac{|Q_{1}|}{r_{1}^{2}} = \frac{8{,}99 \cdot 10^{9} \, \mathrm{\frac{N \cdot m^2}{C^2}} \cdot 7{,}0 \cdot 10^{-6} \, \mathrm{C}}{(0{,}18 \, \mathrm{m})^{2}} = 1{,}9 \cdot 10^{6} \, \mathrm{N/C} \\
                E_{2} &= k_e \frac{|Q_{2}|}{r_{2}^{2}} = \frac{8{,}99 \cdot 10^{9} \, \mathrm{\frac{N \cdot m^2}{C^2}} \cdot 2{,}5 \cdot 10^{-5} \, \mathrm{C}}{(0{,}24 \, \mathrm{m})^{2}} = 3{,}9 \cdot 10^{6} \, \mathrm{N/C}
            \end{align*}
            De rechte hoek in $P$ zorgt ervoor dat $\vec{E}_1$ en $\vec{E}_2$ loodrecht op elkaar staan.
            \begin{align*}
                E_{\text{res}} &= \sqrt{E_1^{2} + E_2^{2}} \simeq 4{,}4 \cdot 10^{6} \, \mathrm{N/C}
            \end{align*}
        \end{solution}
    \end{question}
\end{exercise}
%———————————————————————————————————————

%———————————————————————————————————————
% OEF DRIEHOEK 5: rechthoekige driehoek, cosinusregel
\begin{exercise}
    Twee puntladingen $Q_{A} = -5{,}0 \cdot 10^{-6} \, \mathrm{C}$ en $Q_{B} = -2{,}0 \cdot 10^{-6} \, \mathrm{C}$ bevinden zich in de hoekpunten van een rechthoekige driehoek, zoals te zien op de tekening. De rechte hoek bevindt zich bij $Q_A$. De afstand tussen $Q_A$ en $P$ is $0{,}30 \, \mathrm{m}$, en de afstand tussen $Q_A$ en $Q_B$ is $0{,}40 \, \mathrm{m}$.

    \begin{center}
        \begin{tikzpicture}[scale=0.5]
            \useasboundingbox (-1.5, -1.5) rectangle (5.5, 4.5);
            \tkzDefPoint(0,0){A}
            \begin{scope}[shift=(A)]
                \tkzDefPoint(0:4){B}
                \tkzDefPoint(90:3){C}
            \end{scope}
            \tkzDrawPolygon(A,B,C)

            \tkzLabelPoint[left, below=4pt](A){$Q_{A}$}
            \tkzLabelPoint[right, below=4pt](B){$Q_{B}$}
            \tkzLabelPoint[above=4pt](C){$P$}

            \draw[charge-] (A) circle (\R) node[scale=0.5] {$-$};
            \draw[charge-] (B) circle (\R) node[scale=0.5] {$-$};
            \filldraw (C) circle (2pt);
        \end{tikzpicture}
    \end{center}

    Bereken de grootte van het totale elektrisch veld in punt \textit{P}, veroorzaakt door $Q_A$ en $Q_B$.

    \[ E_{\text{res}} = \answer[tolerance=0.1,onlineshowanswerbutton]{5.5} \cdot 10^{5} \, \mathrm{N/C} \]

    \begin{hint}
        Bereken eerst de afstand $r_{BP}$ met de stelling van Pythagoras in de rechthoekige driehoek.
    \end{hint}
    \begin{hint}
        Bereken de grootte van de twee afzonderlijke velden $E_A$ en $E_B$ in $P$ met de formule van het radiaal veld.
    \end{hint}
    \begin{hint}
        Beide bronladingen zijn negatief: beide velden wijzen dus van $P$ naar hun eigen bronlading toe. De hoek tussen $\vec{E}_A$ en $\vec{E}_B$ is daardoor gelijk aan de hoek van de driehoek in $P$. Bepaal deze hoek met goniometrie ($\tan\theta = \dfrac{\text{overstaande}}{\text{aanliggende}}$), en gebruik vervolgens de cosinusregel met de \textbf{overstaande hoek} in de parallellogram van de veldlijnen om de resultante te vinden.
    \end{hint}

    \begin{solution}\nl
        \underline{Gegevens:}
        \begin{align*}
            & Q_{A} = -5{,}0 \cdot 10^{-6} \, \mathrm{C} & & Q_{B} = -2{,}0 \cdot 10^{-6} \, \mathrm{C} \\
            & r_{AP} = 0{,}30 \, \mathrm{m} & & r_{AB} = 0{,}40 \, \mathrm{m}
        \end{align*}

        \underline{Gevraagd:} $E_{\text{res}}$ in $P$\\

        \underline{Oplossing:}\\
        \textbf{Stap 1: Ontbrekende afstand berekenen}\\
        Met de stelling van Pythagoras in de rechthoekige driehoek ($90^\circ$ bij $Q_A$):
        \begin{align*}
            r_{BP} &= \sqrt{r_{AP}^{2} + r_{AB}^{2}} = \sqrt{(0{,}30 \, \mathrm{m})^{2} + (0{,}40 \, \mathrm{m})^{2}} = 0{,}50 \, \mathrm{m}
        \end{align*}

        \textbf{Stap 2: Grootte berekenen}
        \begin{align*}
            E_{A} &= k_e \frac{|Q_{A}|}{r_{AP}^{2}} = \frac{8{,}99 \cdot 10^{9} \, \mathrm{\frac{N \cdot m^2}{C^2}} \cdot 5{,}0 \cdot 10^{-6} \, \mathrm{C}}{(0{,}30 \, \mathrm{m})^{2}} = 5{,}0 \cdot 10^{5} \, \mathrm{N/C} \\
            E_{B} &= k_e \frac{|Q_{B}|}{r_{BP}^{2}} = \frac{8{,}99 \cdot 10^{9} \, \mathrm{\frac{N \cdot m^2}{C^2}} \cdot 2{,}0 \cdot 10^{-6} \, \mathrm{C}}{(0{,}50 \, \mathrm{m})^{2}} = 7{,}2 \cdot 10^{4} \, \mathrm{N/C}
        \end{align*}

        \textbf{Stap 3: Hoek bepalen}\\
        Omdat $Q_A$ en $Q_B$ beide negatief zijn, wijst $\vec{E}_A$ in $P$ naar $Q_A$ toe, en $\vec{E}_B$ naar $Q_B$ toe. De hoek tussen deze twee veldvectoren is dan gelijk aan de hoek van de driehoek bij $P$:
        \begin{align*}
            \tan(\theta) &= \frac{\text{overstaande}}{\text{aanliggende}} = \frac{r_{AB}}{r_{AP}} = \frac{0{,}40 \, \mathrm{m}}{0{,}30 \, \mathrm{m}} \\
            \theta &= 53{,}1^\circ
        \end{align*}
        De overstaande hoek in het parallellogram is dan $180^\circ - 53{,}1^\circ = 126{,}9^\circ$.

        \textbf{Stap 4: Cosinusregel toepassen}
        \begin{align*}
            E_{\text{res}} &= \sqrt{E_A^{2} + E_B^{2} - 2 \cdot E_A \cdot E_B \cdot \cos(126{,}9^\circ)} \\
                            &= \sqrt{(5{,}0 \cdot 10^{5})^{2} + (7{,}2 \cdot 10^{4})^{2} - 2 \cdot 5{,}0 \cdot 10^{5} \cdot 7{,}2 \cdot 10^{4} \cdot \cos(126{,}9^\circ)} \, \mathrm{N/C} \\
                            &\simeq 5{,}5 \cdot 10^{5} \, \mathrm{N/C}
        \end{align*}
        Het resulterende elektrisch veld in $P$ heeft dus een grootte van ongeveer $5{,}5 \cdot 10^{5} \, \mathrm{N/C}$.
    \end{solution}
\end{exercise}
%———————————————————————————————————————

%———————————————————————————————————————
% OEF DRIEHOEK COS 6: hoek in P gegeven, beide ladingen positief
\begin{exercise}
    Twee positieve puntladingen $Q_{1} = +4{,}0 \cdot 10^{-6} \, \mathrm{C}$ en $Q_{2} = +2{,}0 \cdot 10^{-6} \, \mathrm{C}$ bevinden zich zoals aangegeven op de tekening. De afstand van $Q_1$ tot $P$ is $0{,}25 \, \mathrm{m}$, de afstand van $Q_2$ tot $P$ is $0{,}40 \, \mathrm{m}$, en de hoek tussen beide lijnstukken in $P$ is $70^\circ$.

    \begin{center}
        \begin{tikzpicture}[scale=1]
            \tkzDefPoint(0,0){P}
            \tkzDefPoint(2.5,0){Q1}
            \tkzDefPoint(1.368,3.759){Q2}
            \draw[gray,thick] (P) -- (Q1) node[midway,below]{$0{,}25\,\mathrm{m}$};
            \draw[gray,thick] (P) -- (Q2) node[midway,left]{$0{,}40\,\mathrm{m}$};
            \tkzFillAngle[fill=blue!20,opacity=0.5](Q1,P,Q2)
            \tkzLabelAngle[pos=1.0](Q1,P,Q2){$70^\circ$}
            \tkzLabelPoint[below=1.5mm](P){$P$}
            \tkzLabelPoint[below=1.5mm](Q1){$Q_1$}
            \tkzLabelPoint[above=1.5mm](Q2){$Q_2$}
            \draw[charge+] (Q1) circle (\R) node[scale=0.5]{$+$};
            \draw[charge+] (Q2) circle (\R) node[scale=0.5]{$+$};
            \filldraw (P) circle (1.2pt);
        \end{tikzpicture}
    \end{center}

    Bereken de grootte van het totale elektrisch veld in punt \textit{P}.

    \[ E_{\text{res}} = \answer[tolerance=0.05,onlineshowanswerbutton]{6.2} \cdot 10^{5} \, \mathrm{N/C} \]

    \begin{hint}
        Omdat beide ladingen positief zijn, is de hoek tussen $\vec{E}_1$ en $\vec{E}_2$ gelijk aan de gegeven hoek van $70^\circ$. De overstaande hoek die je in de cosinusregel invult, is dan $180^\circ - 70^\circ$.
    \end{hint}

    \begin{solution}\nl
        \underline{Gegevens:}
        \begin{align*}
            & Q_{1} = +4{,}0 \cdot 10^{-6} \, \mathrm{C} & & Q_{2} = +2{,}0 \cdot 10^{-6} \, \mathrm{C} \\
            & r_{1} = 0{,}25 \, \mathrm{m} & & r_{2} = 0{,}40 \, \mathrm{m}
        \end{align*}

        \underline{Gevraagd:} $E_{\text{res}}$ in $P$\\

        \underline{Oplossing:}\\
        \textbf{Stap 1: Grootte berekenen}
        \begin{align*}
            E_{1} &= k_e \frac{|Q_{1}|}{r_{1}^{2}} = \frac{8{,}99 \cdot 10^{9} \, \mathrm{\frac{N \cdot m^2}{C^2}} \cdot 4{,}0 \cdot 10^{-6} \, \mathrm{C}}{(0{,}25 \, \mathrm{m})^{2}} = 5{,}8 \cdot 10^{5} \, \mathrm{N/C} \\
            E_{2} &= k_e \frac{|Q_{2}|}{r_{2}^{2}} = \frac{8{,}99 \cdot 10^{9} \, \mathrm{\frac{N \cdot m^2}{C^2}} \cdot 2{,}0 \cdot 10^{-6} \, \mathrm{C}}{(0{,}40 \, \mathrm{m})^{2}} = 1{,}1 \cdot 10^{5} \, \mathrm{N/C}
        \end{align*}
        Beide ladingen zijn positief, dus wijzen $\vec{E}_1$ en $\vec{E}_2$ elk weg van hun eigen lading. De hoek tussen deze twee vectoren is dus gelijk aan de gegeven hoek van $70^\circ$. De overstaande hoek in het krachtenparallellogram is $180^\circ - 70^\circ = 110^\circ$.

        \textbf{Stap 2: Cosinusregel toepassen}
        \begin{align*}
            E_{\text{res}} &= \sqrt{E_1^{2} + E_2^{2} - 2 \cdot E_1 \cdot E_2 \cdot \cos(110^\circ)} \\
                            &\simeq 6{,}2 \cdot 10^{5} \, \mathrm{N/C}
        \end{align*}
    \end{solution}
\end{exercise}
%———————————————————————————————————————

%———————————————————————————————————————
% OEF DRIEHOEK COS 7: hoek in P gegeven, tegengesteld teken
\begin{exercise}
    Een positieve puntlading $Q_{1} = +5{,}0 \cdot 10^{-6} \, \mathrm{C}$ en een negatieve puntlading $Q_{3} = -3{,}0 \cdot 10^{-6} \, \mathrm{C}$ bevinden zich zoals aangegeven op de tekening. De afstand van $Q_1$ tot $P$ is $0{,}20 \, \mathrm{m}$, de afstand van $Q_3$ tot $P$ is $0{,}30 \, \mathrm{m}$, en de hoek tussen beide lijnstukken in $P$ is $65^\circ$.

    \begin{center}
        \begin{tikzpicture}[scale=1.3]
            \tkzDefPoint(0,0){P}
            \tkzDefPoint(2.0,0){Q1}
            \tkzDefPoint(1.268,2.719){Q3}
            \draw[gray,thick] (P) -- (Q1) node[midway,below]{$0{,}20\,\mathrm{m}$};
            \draw[gray,thick] (P) -- (Q3) node[midway,left]{$0{,}30\,\mathrm{m}$};
            \tkzFillAngle[fill=blue!20,opacity=0.5](Q1,P,Q3)
            \tkzLabelAngle[pos=1.0](Q1,P,Q3){$65^\circ$}
            \tkzLabelPoint[below=1.5mm](P){$P$}
            \tkzLabelPoint[below=1.5mm](Q1){$Q_1$}
            \tkzLabelPoint[above=1.5mm](Q3){$Q_3$}
            \draw[charge+] (Q1) circle (\R) node[scale=0.5]{$+$};
            \draw[charge-] (Q3) circle (\R) node[scale=0.5]{$-$};
            \filldraw (P) circle (1.2pt);
        \end{tikzpicture}
    \end{center}

    Bereken de grootte van het totale elektrisch veld in punt \textit{P}.

    \[ E_{\text{res}} = \answer[tolerance=0.05,onlineshowanswerbutton]{1.0} \cdot 10^{6} \, \mathrm{N/C} \]

    \begin{hint}
        Schets $\vec{E}_1$ (weg van $Q_1$) en $\vec{E}_3$ (naar $Q_3$ toe). Omdat de ladingen nu een verschillend teken hebben, is de hoek tussen de vectoren zelf niet meer $65^\circ$, maar de overstaande hoek die je uiteindelijk in de cosinusregel gebruikt, blijkt dan net wél gelijk te zijn aan de gegeven hoek van $65^\circ$.
    \end{hint}

    \begin{solution}\nl
        \underline{Gegevens:}
        \begin{align*}
            & Q_{1} = +5{,}0 \cdot 10^{-6} \, \mathrm{C} & & Q_{3} = -3{,}0 \cdot 10^{-6} \, \mathrm{C} \\
            & r_{1} = 0{,}20 \, \mathrm{m} & & r_{3} = 0{,}30 \, \mathrm{m}
        \end{align*}

        \underline{Gevraagd:} $E_{\text{res}}$ in $P$\\

        \underline{Oplossing:}\\
        \textbf{Stap 1: Grootte berekenen}
        \begin{align*}
            E_{1} &= k_e \frac{|Q_{1}|}{r_{1}^{2}} = \frac{8{,}99 \cdot 10^{9} \, \mathrm{\frac{N \cdot m^2}{C^2}} \cdot 5{,}0 \cdot 10^{-6} \, \mathrm{C}}{(0{,}20 \, \mathrm{m})^{2}} = 1{,}1 \cdot 10^{6} \, \mathrm{N/C} \\
            E_{3} &= k_e \frac{|Q_{3}|}{r_{3}^{2}} = \frac{8{,}99 \cdot 10^{9} \, \mathrm{\frac{N \cdot m^2}{C^2}} \cdot 3{,}0 \cdot 10^{-6} \, \mathrm{C}}{(0{,}30 \, \mathrm{m})^{2}} = 3{,}0 \cdot 10^{5} \, \mathrm{N/C}
        \end{align*}
        $\vec{E}_1$ wijst van $P$ weg (langs $PQ_1$), $\vec{E}_3$ wijst van $P$ naar $Q_3$ toe (langs $PQ_3$). Doordat één vector \textit{weg} en de andere \textit{naar toe} wijst, is de overstaande hoek die je in de cosinusregel gebruikt hier gelijk aan de driehoekshoek zelf, dus $65^\circ$ (in plaats van $180^\circ-65^\circ$ zoals bij twee gelijknamige ladingen).

        \textbf{Stap 2: Cosinusregel toepassen}
        \begin{align*}
            E_{\text{res}} &= \sqrt{E_1^{2} + E_3^{2} - 2 \cdot E_1 \cdot E_3 \cdot \cos(65^\circ)} \\
                            &\simeq 1{,}0 \cdot 10^{6} \, \mathrm{N/C}
        \end{align*}
    \end{solution}
\end{exercise}
%———————————————————————————————————————

%———————————————————————————————————————
% OEF DRIEHOEK COS 8: rechthoekige driehoek, rechte hoek NIET in P
\begin{exercise}
    Twee puntladingen $Q_{A} = +6{,}0 \cdot 10^{-6} \, \mathrm{C}$ en $Q_{B} = -3{,}0 \cdot 10^{-6} \, \mathrm{C}$ bevinden zich in de hoekpunten van een rechthoekige driehoek, met de rechte hoek bij $Q_A$. De afstand tussen $Q_A$ en $P$ is $0{,}24 \, \mathrm{m}$, en de afstand tussen $Q_A$ en $Q_B$ is $0{,}32 \, \mathrm{m}$.

    \begin{center}
        \begin{tikzpicture}[scale=0.6]
            \useasboundingbox (-1.5, -1.5) rectangle (5.5, 4.5);
            \tkzDefPoint(0,0){A}
            \begin{scope}[shift=(A)]
                \tkzDefPoint(0:3){P}
                \tkzDefPoint(90:4){B}
            \end{scope}
            \tkzDrawPolygon(A,P,B)
            \tkzMarkRightAngle[size=0.3](P,A,B)

            \tkzLabelPoint[below=1.5mm](A){$Q_{A}$}
            \tkzLabelPoint[below=1.5mm](P){$P$}
            \tkzLabelPoint[above=1.5mm](B){$Q_{B}$}

            \draw[charge+] (A) circle (\R) node[scale=0.5] {$+$};
            \draw[charge-] (B) circle (\R) node[scale=0.5] {$-$};
            \filldraw (P) circle (2pt);
        \end{tikzpicture}
    \end{center}

    Bereken de grootte van het totale elektrisch veld in punt \textit{P}, veroorzaakt door $Q_A$ en $Q_B$.

    \[ E_{\text{res}} = \answer[tolerance=0.1,onlineshowanswerbutton]{8.5} \cdot 10^{5} \, \mathrm{N/C} \]

    \begin{hint}
        Bereken eerst de afstand $r_{BP}$ met de stelling van Pythagoras, en vervolgens de hoek in $P$ met $\tan\theta = \dfrac{\text{overstaande}}{\text{aanliggende}}$.
    \end{hint}
    \begin{hint}
        $Q_A$ is positief en $Q_B$ is negatief, net als in een vorige oefening betekent dit dat de overstaande hoek in de cosinusregel gelijk is aan de hoek van de driehoek zelf.
    \end{hint}

    \begin{solution}\nl
        \underline{Gegevens:}
        \begin{align*}
            & Q_{A} = +6{,}0 \cdot 10^{-6} \, \mathrm{C} & & Q_{B} = -3{,}0 \cdot 10^{-6} \, \mathrm{C} \\
            & r_{AP} = 0{,}24 \, \mathrm{m} & & r_{AB} = 0{,}32 \, \mathrm{m}
        \end{align*}

        \underline{Gevraagd:} $E_{\text{res}}$ in $P$\\

        \underline{Oplossing:}\\
        \textbf{Stap 1: Ontbrekende afstand}\\
        Met de stelling van Pythagoras in de rechthoekige driehoek ($90^\circ$ bij $Q_A$):
        \begin{align*}
            r_{BP} &= \sqrt{r_{AP}^{2} + r_{AB}^{2}} = \sqrt{(0{,}24 \, \mathrm{m})^{2} + (0{,}32 \, \mathrm{m})^{2}} = 0{,}40 \, \mathrm{m}
        \end{align*}

        \textbf{Stap 2: Grootte berekenen}
        \begin{align*}
            E_{A} &= k_e \frac{|Q_{A}|}{r_{AP}^{2}} = \frac{8{,}99 \cdot 10^{9} \, \mathrm{\frac{N \cdot m^2}{C^2}} \cdot 6{,}0 \cdot 10^{-6} \, \mathrm{C}}{(0{,}24 \, \mathrm{m})^{2}} = 9{,}4 \cdot 10^{5} \, \mathrm{N/C} \\
            E_{B} &= k_e \frac{|Q_{B}|}{r_{BP}^{2}} = \frac{8{,}99 \cdot 10^{9} \, \mathrm{\frac{N \cdot m^2}{C^2}} \cdot 3{,}0 \cdot 10^{-6} \, \mathrm{C}}{(0{,}40 \, \mathrm{m})^{2}} = 1{,}7 \cdot 10^{5} \, \mathrm{N/C}
        \end{align*}

        \textbf{Stap 3: Hoek bepalen}
        \begin{align*}
            \tan(\theta) &= \frac{r_{AB}}{r_{AP}} = \frac{0{,}32 \, \mathrm{m}}{0{,}24 \, \mathrm{m}} \quad\Rightarrow\quad \theta = 53{,}1^\circ
        \end{align*}
        $Q_A$ is positief ($\vec{E}_A$ wijst weg van $Q_A$) en $Q_B$ is negatief ($\vec{E}_B$ wijst naar $Q_B$ toe): de overstaande hoek in de cosinusregel is dan gelijk aan de driehoekshoek zelf, $53{,}1^\circ$.

        \textbf{Stap 4: Cosinusregel toepassen}
        \begin{align*}
            E_{\text{res}} &= \sqrt{E_A^{2} + E_B^{2} - 2 \cdot E_A \cdot E_B \cdot \cos(53{,}1^\circ)} \\
                            &\simeq 8{,}5 \cdot 10^{5} \, \mathrm{N/C}
        \end{align*}
    \end{solution}
\end{exercise}
%———————————————————————————————————————

%———————————————————————————————————————
% OEF DRIEHOEK COS 9: willekeurige driehoek, hoek via cosinusregel uit 3 zijden
\begin{exercise}
    Twee puntladingen $Q_{1} = +7{,}0 \cdot 10^{-6} \, \mathrm{C}$ en $Q_{2} = -5{,}0 \cdot 10^{-6} \, \mathrm{C}$ bevinden zich in twee hoekpunten van een driehoek met punt $P$, zoals aangegeven op de tekening. De afstand van $Q_1$ tot $P$ is $0{,}15 \, \mathrm{m}$, de afstand van $Q_2$ tot $P$ is $0{,}20 \, \mathrm{m}$, en de afstand tussen $Q_1$ en $Q_2$ onderling is $0{,}30 \, \mathrm{m}$.

    \begin{center}
        \begin{tikzpicture}[scale=1.3]
            \tkzDefPoint(0,0){P}
            \tkzDefPoint(1.5,0){Q1}
            \tkzDefPoint(-0.917,1.778){Q2}
            \tkzDrawPolygon(P,Q1,Q2)
            \tkzLabelSegment[below=1.5mm](P,Q1){$0{,}15\,\mathrm{m}$}
            \tkzLabelSegment[above=1.5mm](P,Q2){$0{,}20\,\mathrm{m}$}
            \tkzLabelSegment[right=1.5mm](Q1,Q2){$0{,}30\,\mathrm{m}$}
            \tkzLabelPoint[below=1.5mm](P){$P$}
            \tkzLabelPoint[right=1.5mm](Q1){$Q_1$}
            \tkzLabelPoint[above=1.5mm](Q2){$Q_2$}
            \draw[charge+] (Q1) circle (\R) node[scale=0.5]{$+$};
            \draw[charge-] (Q2) circle (\R) node[scale=0.5]{$-$};
            \filldraw (P) circle (1.2pt);
        \end{tikzpicture}
    \end{center}

    Bereken de grootte van het totale elektrisch veld in punt \textit{P}.

    \[ E_{\text{res}} = \answer[tolerance=0.1,onlineshowanswerbutton]{3.5} \cdot 10^{6} \, \mathrm{N/C} \]

    \begin{hint}
        Deze keer is geen enkele hoek rechtstreeks gegeven, en de driehoek is niet rechthoekig. Bereken eerst de hoek in $P$ met de cosinusregel, toegepast op de \textbf{driehoek van de afstanden} (met de drie gegeven zijden).
    \end{hint}
    \begin{hint}
        Zodra je de hoek in $P$ kent, kan je op dezelfde manier als in de vorige oefeningen verder werken om de overstaande hoek en vervolgens $E_{\text{res}}$ te bepalen.
    \end{hint}

    \begin{solution}\nl
        \underline{Gegevens:}
        \begin{align*}
            & Q_{1} = +7{,}0 \cdot 10^{-6} \, \mathrm{C} & & Q_{2} = -5{,}0 \cdot 10^{-6} \, \mathrm{C} \\
            & r_{1} = 0{,}15 \, \mathrm{m} & & r_{2} = 0{,}20 \, \mathrm{m} & & r_{12} = 0{,}30 \, \mathrm{m}
        \end{align*}

        \underline{Gevraagd:} $E_{\text{res}}$ in $P$\\

        \underline{Oplossing:}\\
        \textbf{Stap 1: Hoek in $P$ bepalen met de cosinusregel}\\
        We passen de cosinusregel toe op de driehoek van de afstanden, met $r_{12}$ als de zijde tegenover de hoek in $P$:
        \begin{align*}
            r_{12}^{2} &= r_{1}^{2} + r_{2}^{2} - 2 \cdot r_{1} \cdot r_{2} \cdot \cos(P) \\
            \cos(P) &= \frac{r_{1}^{2} + r_{2}^{2} - r_{12}^{2}}{2 \cdot r_{1} \cdot r_{2}} = \frac{(0{,}15)^{2} + (0{,}20)^{2} - (0{,}30)^{2}}{2 \cdot 0{,}15 \cdot 0{,}20} = -0{,}458 \\
            P &= 117{,}3^\circ
        \end{align*}

        \textbf{Stap 2: Grootte van de velden berekenen}
        \begin{align*}
            E_{1} &= k_e \frac{|Q_{1}|}{r_{1}^{2}} = \frac{8{,}99 \cdot 10^{9} \, \mathrm{\frac{N \cdot m^2}{C^2}} \cdot 7{,}0 \cdot 10^{-6} \, \mathrm{C}}{(0{,}15 \, \mathrm{m})^{2}} = 2{,}8 \cdot 10^{6} \, \mathrm{N/C} \\
            E_{2} &= k_e \frac{|Q_{2}|}{r_{2}^{2}} = \frac{8{,}99 \cdot 10^{9} \, \mathrm{\frac{N \cdot m^2}{C^2}} \cdot 5{,}0 \cdot 10^{-6} \, \mathrm{C}}{(0{,}20 \, \mathrm{m})^{2}} = 1{,}1 \cdot 10^{6} \, \mathrm{N/C}
        \end{align*}

        \textbf{Stap 3: Overstaande hoek en cosinusregel}\\
        $Q_1$ is positief ($\vec{E}_1$ wijst weg van $Q_1$) en $Q_2$ is negatief ($\vec{E}_2$ wijst naar $Q_2$ toe): net als in de vorige oefeningen is de overstaande hoek in de cosinusregel dan gelijk aan de hoek $P$ zelf, dus $117{,}3^\circ$.
        \begin{align*}
            E_{\text{res}} &= \sqrt{E_1^{2} + E_2^{2} - 2 \cdot E_1 \cdot E_2 \cdot \cos(117{,}3^\circ)} \\
                            &\simeq 3{,}5 \cdot 10^{6} \, \mathrm{N/C}
        \end{align*}
    \end{solution}
\end{exercise}
%———————————————————————————————————————

%———————————————————————————————————————
% OEF VELD COS 6: hoek gegeven, beide ladingen positief (nC en µC)
\begin{exercise}
    Twee positieve puntladingen $Q_{1} = +600 \, \mathrm{nC}$ en $Q_{2} = +1{,}5 \, \mathrm{\mu C}$ bevinden zich zoals aangegeven op de tekening. De afstand van $Q_1$ tot $P$ is $12 \, \mathrm{cm}$, de afstand van $Q_2$ tot $P$ is $18 \, \mathrm{cm}$, en de hoek tussen beide lijnstukken in $P$ is $60^\circ$.

    \begin{center}
        \begin{tikzpicture}[scale=1.3]
            \tkzDefPoint(0,0){P}
            \tkzDefPoint(1.2,0){Q1}
            \tkzDefPoint(0.9,1.559){Q2}
            \draw[gray,thick] (P) -- (Q1) node[midway,below]{$12\,\mathrm{cm}$};
            \draw[gray,thick] (P) -- (Q2) node[midway,left]{$18\,\mathrm{cm}$};
            \tkzFillAngle[fill=blue!20,opacity=0.5](Q1,P,Q2)
            \tkzLabelAngle[pos=1.0](Q1,P,Q2){$60^\circ$}
            \tkzLabelPoint[below=1.5mm](P){$P$}
            \tkzLabelPoint[below=1.5mm](Q1){$Q_1$}
            \tkzLabelPoint[above=1.5mm](Q2){$Q_2$}
            \draw[charge+] (Q1) circle (\R) node[scale=0.5]{$+$};
            \draw[charge+] (Q2) circle (\R) node[scale=0.5]{$+$};
            \filldraw (P) circle (1.2pt);
        \end{tikzpicture}
    \end{center}

    Bereken de grootte van het totale elektrisch veld in punt \textit{P}.

    \[ E_{\text{res}} = \answer[tolerance=0.05,onlineshowanswerbutton]{6.9} \cdot 10^{5} \, \mathrm{N/C} \]

    \begin{hint}
        Beide ladingen zijn positief: de overstaande hoek in de cosinusregel is $180^\circ - 60^\circ$.
    \end{hint}

    \begin{solution}\nl
        \underline{Gegevens:}
        \begin{align*}
            & Q_{1} = +600 \, \mathrm{nC} = 6{,}0 \cdot 10^{-7} \, \mathrm{C} & & Q_{2} = +1{,}5 \, \mathrm{\mu C} = 1{,}5 \cdot 10^{-6} \, \mathrm{C} \\
            & r_{1} = 12 \, \mathrm{cm} = 0{,}12 \, \mathrm{m} & & r_{2} = 18 \, \mathrm{cm} = 0{,}18 \, \mathrm{m}
        \end{align*}

        \underline{Gevraagd:} $E_{\text{res}}$ in $P$\\

        \underline{Oplossing:}\\
        \begin{align*}
            E_{1} &= k_e \frac{|Q_{1}|}{r_{1}^{2}} = \frac{8{,}99 \cdot 10^{9} \, \mathrm{\frac{N \cdot m^2}{C^2}} \cdot 6{,}0 \cdot 10^{-7} \, \mathrm{C}}{(0{,}12 \, \mathrm{m})^{2}} = 3{,}7 \cdot 10^{5} \, \mathrm{N/C} \\
            E_{2} &= k_e \frac{|Q_{2}|}{r_{2}^{2}} = \frac{8{,}99 \cdot 10^{9} \, \mathrm{\frac{N \cdot m^2}{C^2}} \cdot 1{,}5 \cdot 10^{-6} \, \mathrm{C}}{(0{,}18 \, \mathrm{m})^{2}} = 4{,}2 \cdot 10^{5} \, \mathrm{N/C}
        \end{align*}
        Beide ladingen zijn positief, dus is de hoek tussen $\vec{E}_1$ en $\vec{E}_2$ gelijk aan $60^\circ$. De overstaande hoek in de cosinusregel is $180^\circ - 60^\circ = 120^\circ$.
        \begin{align*}
            E_{\text{res}} &= \sqrt{E_1^{2} + E_2^{2} - 2 \cdot E_1 \cdot E_2 \cdot \cos(120^\circ)} \simeq 6{,}9 \cdot 10^{5} \, \mathrm{N/C}
        \end{align*}
    \end{solution}
\end{exercise}
%———————————————————————————————————————

%———————————————————————————————————————
% OEF VELD COS 7: hoek gegeven, tegengesteld teken (nC, mm)
\begin{exercise}
    Een positieve puntlading $Q_{1} = +40 \, \mathrm{nC}$ en een negatieve puntlading $Q_{3} = -90 \, \mathrm{nC}$ bevinden zich zoals aangegeven op de tekening. De afstand van $Q_1$ tot $P$ is $60 \, \mathrm{mm}$, de afstand van $Q_3$ tot $P$ is $90 \, \mathrm{mm}$, en de hoek tussen beide lijnstukken in $P$ is $75^\circ$.

    \begin{center}
        \begin{tikzpicture}[scale=1.6]
            \tkzDefPoint(0,0){P}
            \tkzDefPoint(0.6,0){Q1}
            \tkzDefPoint(0.233,0.869){Q3}
            \draw[gray,thick] (P) -- (Q1) node[midway,below]{$60\,\mathrm{mm}$};
            \draw[gray,thick] (P) -- (Q3) node[midway,left]{$90\,\mathrm{mm}$};
            \tkzFillAngle[fill=blue!20,opacity=0.5](Q1,P,Q3)
            \tkzLabelAngle[pos=1.0](Q1,P,Q3){$75^\circ$}
            \tkzLabelPoint[below=1.5mm](P){$P$}
            \tkzLabelPoint[below=1.5mm](Q1){$Q_1$}
            \tkzLabelPoint[above=1.5mm](Q3){$Q_3$}
            \draw[charge+] (Q1) circle (\R) node[scale=0.5]{$+$};
            \draw[charge-] (Q3) circle (\R) node[scale=0.5]{$-$};
            \filldraw (P) circle (1.2pt);
        \end{tikzpicture}
    \end{center}

    Bereken de grootte van het totale elektrisch veld in punt \textit{P}.

    \[ E_{\text{res}} = \answer[tolerance=0.01,onlineshowanswerbutton]{1.2} \cdot 10^{5} \, \mathrm{N/C} \]

    \begin{hint}
        Verschillend teken: de overstaande hoek in de cosinusregel is hier gelijk aan de gegeven hoek zelf, dus $75^\circ$.
    \end{hint}

    \begin{solution}\nl
        \underline{Gegevens:}
        \begin{align*}
            & Q_{1} = +40 \, \mathrm{nC} = 4{,}0 \cdot 10^{-8} \, \mathrm{C} & & Q_{3} = -90 \, \mathrm{nC} = 9{,}0 \cdot 10^{-8} \, \mathrm{C} \\
            & r_{1} = 60 \, \mathrm{mm} = 0{,}060 \, \mathrm{m} & & r_{3} = 90 \, \mathrm{mm} = 0{,}090 \, \mathrm{m}
        \end{align*}

        \underline{Gevraagd:} $E_{\text{res}}$ in $P$\\

        \underline{Oplossing:}\\
        \begin{align*}
            E_{1} &= k_e \frac{|Q_{1}|}{r_{1}^{2}} = \frac{8{,}99 \cdot 10^{9} \, \mathrm{\frac{N \cdot m^2}{C^2}} \cdot 4{,}0 \cdot 10^{-8} \, \mathrm{C}}{(0{,}060 \, \mathrm{m})^{2}} = 1{,}0 \cdot 10^{5} \, \mathrm{N/C} \\
            E_{3} &= k_e \frac{|Q_{3}|}{r_{3}^{2}} = \frac{8{,}99 \cdot 10^{9} \, \mathrm{\frac{N \cdot m^2}{C^2}} \cdot 9{,}0 \cdot 10^{-8} \, \mathrm{C}}{(0{,}090 \, \mathrm{m})^{2}} = 1{,}0 \cdot 10^{5} \, \mathrm{N/C}
        \end{align*}
        Doordat $Q_1$ en $Q_3$ een verschillend teken hebben, is de overstaande hoek in de cosinusregel gelijk aan de gegeven hoek zelf: $75^\circ$.
        \begin{align*}
            E_{\text{res}} &= \sqrt{E_1^{2} + E_3^{2} - 2 \cdot E_1 \cdot E_3 \cdot \cos(75^\circ)} \simeq 1{,}2 \cdot 10^{5} \, \mathrm{N/C}
        \end{align*}
    \end{solution}
\end{exercise}
%———————————————————————————————————————

%———————————————————————————————————————
% OEF VELD COS 8: rechthoekige driehoek, niet in P, beide ladingen positief
\begin{exercise}
    Twee positieve puntladingen $Q_{A} = +3{,}0 \, \mathrm{\mu C}$ en $Q_{B} = +4{,}0 \, \mathrm{\mu C}$ bevinden zich in de hoekpunten van een rechthoekige driehoek, met de rechte hoek bij $Q_A$. De afstand tussen $Q_A$ en $P$ is $9{,}0 \, \mathrm{cm}$, en de afstand tussen $Q_A$ en $Q_B$ is $12 \, \mathrm{cm}$.

    \begin{center}
        \begin{tikzpicture}[scale=0.5]
            \useasboundingbox (-1.5, -1.5) rectangle (5.5, 4.5);
            \tkzDefPoint(0,0){A}
            \begin{scope}[shift=(A)]
                \tkzDefPoint(0:3){P}
                \tkzDefPoint(90:4){B}
            \end{scope}
            \tkzDrawPolygon(A,P,B)
            \tkzMarkRightAngle[size=0.3](P,A,B)
            \tkzLabelPoint[below=1.5mm](A){$Q_{A}$}
            \tkzLabelPoint[below=1.5mm](P){$P$}
            \tkzLabelPoint[above=1.5mm](B){$Q_{B}$}
            \draw[charge+] (A) circle (\R) node[scale=0.5] {$+$};
            \draw[charge+] (B) circle (\R) node[scale=0.5] {$+$};
            \filldraw (P) circle (2pt);
        \end{tikzpicture}
    \end{center}

    Bereken de grootte van het totale elektrisch veld in punt \textit{P}.

    \[ E_{\text{res}} = \answer[tolerance=0.1,onlineshowanswerbutton]{4.5} \cdot 10^{6} \, \mathrm{N/C} \]

    \begin{hint}
        Bereken eerst $r_{BP}$ met Pythagoras en de hoek in $P$ met $\tan\theta$. Let op: dit keer zijn beide ladingen positief.
    \end{hint}

    \begin{solution}\nl
        \underline{Gegevens:}
        \begin{align*}
            & Q_{A} = +3{,}0 \, \mathrm{\mu C} = 3{,}0 \cdot 10^{-6} \, \mathrm{C} & & Q_{B} = +4{,}0 \, \mathrm{\mu C} = 4{,}0 \cdot 10^{-6} \, \mathrm{C} \\
            & r_{AP} = 9{,}0 \, \mathrm{cm} = 0{,}090 \, \mathrm{m} & & r_{AB} = 12 \, \mathrm{cm} = 0{,}12 \, \mathrm{m}
        \end{align*}

        \underline{Gevraagd:} $E_{\text{res}}$ in $P$\\

        \underline{Oplossing:}\\
        \textbf{Stap 1: Ontbrekende afstand}
        \begin{align*}
            r_{BP} &= \sqrt{r_{AP}^{2} + r_{AB}^{2}} = \sqrt{(0{,}090 \, \mathrm{m})^{2} + (0{,}12 \, \mathrm{m})^{2}} = 0{,}15 \, \mathrm{m}
        \end{align*}

        \textbf{Stap 2: Grootte berekenen}
        \begin{align*}
            E_{A} &= k_e \frac{|Q_{A}|}{r_{AP}^{2}} = \frac{8{,}99 \cdot 10^{9} \, \mathrm{\frac{N \cdot m^2}{C^2}} \cdot 3{,}0 \cdot 10^{-6} \, \mathrm{C}}{(0{,}090 \, \mathrm{m})^{2}} = 3{,}3 \cdot 10^{6} \, \mathrm{N/C} \\
            E_{B} &= k_e \frac{|Q_{B}|}{r_{BP}^{2}} = \frac{8{,}99 \cdot 10^{9} \, \mathrm{\frac{N \cdot m^2}{C^2}} \cdot 4{,}0 \cdot 10^{-6} \, \mathrm{C}}{(0{,}15 \, \mathrm{m})^{2}} = 1{,}6 \cdot 10^{6} \, \mathrm{N/C}
        \end{align*}

        \textbf{Stap 3: Hoek en cosinusregel}\\
        $\tan\theta = \dfrac{r_{AB}}{r_{AP}} = \dfrac{0{,}12}{0{,}090} \Rightarrow \theta = 53{,}1^\circ$. Beide ladingen zijn nu positief, dus is de overstaande hoek $180^\circ - 53{,}1^\circ = 126{,}9^\circ$.
        \begin{align*}
            E_{\text{res}} &= \sqrt{E_A^{2} + E_B^{2} - 2 \cdot E_A \cdot E_B \cdot \cos(126{,}9^\circ)} \simeq 4{,}5 \cdot 10^{6} \, \mathrm{N/C}
        \end{align*}
    \end{solution}
\end{exercise}
%———————————————————————————————————————

%———————————————————————————————————————
% OEF VELD COS 9: rechthoekige driehoek niet in P, beide negatief, grote onderlinge afstand
\begin{exercise}
    Twee negatieve puntladingen $Q_{A} = -8{,}0 \, \mathrm{\mu C}$ en $Q_{B} = -5{,}0 \, \mathrm{\mu C}$ bevinden zich in de hoekpunten van een rechthoekige driehoek, met de rechte hoek bij $Q_A$. De afstand tussen $Q_A$ en $P$ is $14 \, \mathrm{cm}$, en de afstand tussen $Q_A$ en $Q_B$ zelf is $50 \, \mathrm{cm}$.

    \begin{center}
        \begin{tikzpicture}[scale=0.15]
            \useasboundingbox (-5, -5) rectangle (55, 55);
            \tkzDefPoint(0,0){A}
            \begin{scope}[shift=(A)]
                \tkzDefPoint(0:14){P}
                \tkzDefPoint(90:48){B}
            \end{scope}
            \tkzDrawPolygon(A,P,B)
            \tkzMarkRightAngle[size=4](P,A,B)
            \tkzLabelPoint[below=1.5mm](A){$Q_{A}$}
            \tkzLabelPoint[below=1.5mm](P){$P$}
            \tkzLabelPoint[above=1.5mm](B){$Q_{B}$}
            \draw[charge-] (A) circle (2.2) node[scale=1.3] {$-$};
            \draw[charge-] (B) circle (2.2) node[scale=1.3] {$-$};
            \filldraw (P) circle (0.8);
        \end{tikzpicture}
    \end{center}

    Bereken de grootte van het totale elektrisch veld in punt \textit{P}.

    \[ E_{\text{res}} = \answer[tolerance=0.1,onlineshowanswerbutton]{3.7} \cdot 10^{6} \, \mathrm{N/C} \]

    \begin{hint}
        De onderlinge afstand van $50 \, \mathrm{cm}$ is de schuine zijde van de rechthoekige driehoek. Bereken $r_{BP}$ met Pythagoras.
    \end{hint}
    \begin{hint}
        Beide ladingen zijn negatief: dat is dezelfde situatie als bij twee positieve ladingen (overstaande hoek $=180^\circ - \theta$).
    \end{hint}

    \begin{solution}\nl
        \underline{Gegevens:}
        \begin{align*}
            & Q_{A} = -8{,}0 \, \mathrm{\mu C} & & Q_{B} = -5{,}0 \, \mathrm{\mu C} \\
            & r_{AP} = 14 \, \mathrm{cm} = 0{,}14 \, \mathrm{m} & & r_{AB} = 50 \, \mathrm{cm} = 0{,}50 \, \mathrm{m}
        \end{align*}

        \underline{Gevraagd:} $E_{\text{res}}$ in $P$\\

        \underline{Oplossing:}\\
        \textbf{Stap 1: Ontbrekende afstand}
        \begin{align*}
            r_{BP} &= \sqrt{r_{AB}^{2} - r_{AP}^{2}} = \sqrt{(0{,}50 \, \mathrm{m})^{2} - (0{,}14 \, \mathrm{m})^{2}} = 0{,}48 \, \mathrm{m}
        \end{align*}

        \textbf{Stap 2: Grootte berekenen}
        \begin{align*}
            E_{A} &= k_e \frac{|Q_{A}|}{r_{AP}^{2}} = \frac{8{,}99 \cdot 10^{9} \, \mathrm{\frac{N \cdot m^2}{C^2}} \cdot 8{,}0 \cdot 10^{-6} \, \mathrm{C}}{(0{,}14 \, \mathrm{m})^{2}} = 3{,}7 \cdot 10^{6} \, \mathrm{N/C} \\
            E_{B} &= k_e \frac{|Q_{B}|}{r_{BP}^{2}} = \frac{8{,}99 \cdot 10^{9} \, \mathrm{\frac{N \cdot m^2}{C^2}} \cdot 5{,}0 \cdot 10^{-6} \, \mathrm{C}}{(0{,}48 \, \mathrm{m})^{2}} = 2{,}0 \cdot 10^{5} \, \mathrm{N/C}
        \end{align*}

        \textbf{Stap 3: Hoek en cosinusregel}\\
        $\tan\theta = \dfrac{r_{BP}}{r_{AP}} = \dfrac{0{,}48}{0{,}14} \Rightarrow \theta = 73{,}7^\circ$. Beide ladingen zijn negatief (zelfde teken), dus is de overstaande hoek $180^\circ - 73{,}7^\circ = 106{,}3^\circ$.
        \begin{align*}
            E_{\text{res}} &= \sqrt{E_A^{2} + E_B^{2} - 2 \cdot E_A \cdot E_B \cdot \cos(106{,}3^\circ)} \simeq 3{,}7 \cdot 10^{6} \, \mathrm{N/C}
        \end{align*}
        Merk op dat $E_{\text{res}}$ hier nauwelijks verschilt van $E_A$: omdat $Q_A$ zoveel dichter bij $P$ ligt, domineert die bijdrage volledig.
    \end{solution}
\end{exercise}
%———————————————————————————————————————

%———————————————————————————————————————
% OEF VELD COS 10: willekeurige driehoek, hoek via cosinusregel, beide positief
\begin{exercise}
    Twee positieve puntladingen $Q_{1} = +6{,}0 \, \mathrm{\mu C}$ en $Q_{2} = +2{,}5 \, \mathrm{\mu C}$ bevinden zich in twee hoekpunten van een driehoek met punt $P$. De afstand van $Q_1$ tot $P$ is $10 \, \mathrm{cm}$, de afstand van $Q_2$ tot $P$ is $14 \, \mathrm{cm}$, en de afstand tussen $Q_1$ en $Q_2$ onderling is $20 \, \mathrm{cm}$.

    \begin{center}
        \begin{tikzpicture}[scale=1.3]
            \tkzDefPoint(0,0){P}
            \tkzDefPoint(1.0,0){Q1}
            \tkzDefPoint(-0.520,1.325){Q2}
            \tkzDrawPolygon(P,Q1,Q2)
            \tkzLabelSegment[below=1.5mm](P,Q1){$10\,\mathrm{cm}$}
            \tkzLabelSegment[above=1.5mm](P,Q2){$14\,\mathrm{cm}$}
            \tkzLabelSegment[right=1.5mm](Q1,Q2){$20\,\mathrm{cm}$}
            \tkzLabelPoint[below=1.5mm](P){$P$}
            \tkzLabelPoint[right=1.5mm](Q1){$Q_1$}
            \tkzLabelPoint[above=1.5mm](Q2){$Q_2$}
            \draw[charge+] (Q1) circle (\R) node[scale=0.5]{$+$};
            \draw[charge+] (Q2) circle (\R) node[scale=0.5]{$+$};
            \filldraw (P) circle (1.2pt);
        \end{tikzpicture}
    \end{center}

    Bereken de grootte van het totale elektrisch veld in punt \textit{P}.

    \[ E_{\text{res}} = \answer[tolerance=0.1,onlineshowanswerbutton]{5.1} \cdot 10^{6} \, \mathrm{N/C} \]

    \begin{hint}
        Bereken eerst de hoek in $P$ met de cosinusregel, toegepast op de driehoek van de afstanden.
    \end{hint}

    \begin{solution}\nl
        \underline{Gegevens:}
        \begin{align*}
            & Q_{1} = +6{,}0 \, \mathrm{\mu C} & & Q_{2} = +2{,}5 \, \mathrm{\mu C} \\
            & r_{1} = 10 \, \mathrm{cm} = 0{,}10 \, \mathrm{m} & & r_{2} = 14 \, \mathrm{cm} = 0{,}14 \, \mathrm{m} & & r_{12} = 20 \, \mathrm{cm} = 0{,}20 \, \mathrm{m}
        \end{align*}

        \underline{Gevraagd:} $E_{\text{res}}$ in $P$\\

        \underline{Oplossing:}\\
        \textbf{Stap 1: Hoek in $P$}
        \begin{align*}
            \cos(P) &= \frac{r_{1}^{2} + r_{2}^{2} - r_{12}^{2}}{2 \cdot r_{1} \cdot r_{2}} = \frac{(0{,}10)^{2} + (0{,}14)^{2} - (0{,}20)^{2}}{2 \cdot 0{,}10 \cdot 0{,}14} = -0{,}371 \\
            P &= 111{,}8^\circ
        \end{align*}

        \textbf{Stap 2: Grootte berekenen}
        \begin{align*}
            E_{1} &= k_e \frac{|Q_{1}|}{r_{1}^{2}} = \frac{8{,}99 \cdot 10^{9} \, \mathrm{\frac{N \cdot m^2}{C^2}} \cdot 6{,}0 \cdot 10^{-6} \, \mathrm{C}}{(0{,}10 \, \mathrm{m})^{2}} = 5{,}4 \cdot 10^{6} \, \mathrm{N/C} \\
            E_{2} &= k_e \frac{|Q_{2}|}{r_{2}^{2}} = \frac{8{,}99 \cdot 10^{9} \, \mathrm{\frac{N \cdot m^2}{C^2}} \cdot 2{,}5 \cdot 10^{-6} \, \mathrm{C}}{(0{,}14 \, \mathrm{m})^{2}} = 1{,}1 \cdot 10^{6} \, \mathrm{N/C}
        \end{align*}

        \textbf{Stap 3: Overstaande hoek en cosinusregel}\\
        Beide ladingen zijn positief (zelfde teken), dus is de overstaande hoek $180^\circ - 111{,}8^\circ = 68{,}2^\circ$.
        \begin{align*}
            E_{\text{res}} &= \sqrt{E_1^{2} + E_2^{2} - 2 \cdot E_1 \cdot E_2 \cdot \cos(68{,}2^\circ)} \simeq 5{,}1 \cdot 10^{6} \, \mathrm{N/C}
        \end{align*}
    \end{solution}
\end{exercise}
%———————————————————————————————————————

\end{document}